{\large\textbf{Solenoid and loop}}

A closed circular loop of radius $r$ consists of an ideal battery of
electromotive force $\xi$ and a wire of resistance $R$. A long thin air-core
solenoid is aligned with the axis of the loop ($z$-axis). Its length is
$l \gg r$, cross-sectional area is $A \left(r \gg \sqrt{A}\right)$, and the
number of turns is $N$. The solenoid is powered by a constant current $I$
provided by an ideal current source. The directions of the currents in the
solenoid and in the loop are the same (clockwise in the figure).

\begin{center}\includegraphics[width=0.67\linewidth]{figures/EuPhO_2020_Q1_fig1.png}\end{center}

a) Find the force $F_1$ acting on the solenoid when its head $O_1$ is
positioned in the loop centre $O$. What is the force $F_2$ acting on the
solenoid when its tail $O_2$ is located in the centre of the loop?

b) Suppose now, that the solenoid is moving slowly with a constant velocity $v$
along $z$-axis starting far away from the loop, going past its centre, and
proceeding further to the right in positive $z$-direction. Plot the current $J$
flowing in the loop as a function of time. Highlight important features and
values on the graph. The velocity $v$ is so small that self inductance of the
loop can be neglected.
